\section{The polynomial that isolates one edge}\label{sec:symbols}

\status{The localizer, its elementary keys, the finite correction process, and the degree accounting are proved here. Passing from the selected path-polynomial space to the key-symbol algebra remains an explicit representation-theoretic input. Thus this section opens, but does not by itself finish, the Type-II construction in JR9.} \textsf{[SY1--SY8]}

Work in the Type-II setting: $p\ge13$, $2\le r<k$, an even Frobenius period $n'$, and a coefficient DVR $O$ with uniformizer $w$ and $w^2=p$. A frame consists of disjoint increasing blocks of even positions $a_{j,1}<\cdots<a_{j,r}=z_j$, periodic under $j\mapsto j+J$ by addition of $n'$. Put $Q_j=p^{z_{j+1}-z_j}$ and let $Q=p^{G'}$ bound all Frobenius exponents used below. The full frame and size conditions are those in the ledger; in particular, they are stronger than just periodicity.

The chart substitution $\Phi_g$ is an integral invertible change of coordinates followed by multiplication of the normal variables by $w$; let $\Pi(g)$ be its rational plane. At an even position the normal variables have column $b\ge r$; at an odd position they have $b<r$. Its inverse exists over $\operatorname{Frac}(O)$. Write $\mathcal T$ for the integral polynomials in the remaining, tangential variables. All variables here have positions in the selected set $TT$.

\subsection{A finite-field telescoping identity}

Let $M_{j,I}$ be the $r$ by $r$ minor, in columns $I$, of the rows obtained by raising row $a_{j,l}$ to the power $p^{z_j-a_{j,l}}$. Let $I_0=\{0,\ldots,r-1\}$, $\Delta_j=M_{j,I_0}$, and let $\mathcal I$ consist of the $r(k-r)$ sets obtained by replacing one column of $I_0$. Define
\begin{align}
R_{\rm loc}&=\Delta_0\prod_{I\in\mathcal I}
 \left(\prod_{j=0}^{J-1}\Delta_j^{Q_j-1}
             -\prod_{j=0}^{J-1}M_{j,I}^{Q_j-1}\right),\label{eq:sy-localizer}\\
R_0&=\prod_{j=0}^{J-1}\Delta_j^{\alpha_j},\qquad
\alpha_j=\mathbf1_{j=0}+r(k-r)(Q_j-1).
\end{align}
These are exactly the source's \texttt{Rloc} and \texttt{R0}.\cite{symbol-source,symbol-elementary}

\begin{lemma}[Plane isolation]\label{lem:sy-isolation}
For every other rational $r$-plane chart $g'$, all coefficients of $\Phi_{g'}R_{\rm loc}$ are divisible by $w$. In its own chart, $\Phi_gR_{\rm loc}\equiv\Phi_gR_0\pmod w$.
\end{lemma}
\begin{proof}
Translate all relevant lifted positions by a common multiple of $n'$ to make them nonnegative; this changes neither differences nor coordinates. Reduce modulo $w$, and put $q_0=p^{n'}$. In the other chart the rows, after their indicated Frobenius powers, have the form $T_jH^{(p^{z_j})}$, with $H\in\operatorname{Mat}_{r,k}(\mathbf F_{q_0})$. Hence
\[
M_{j,I}=\det(T_j)\det(H_I)^{p^{z_j}},\qquad
\sum_j p^{z_j}(Q_j-1)=p^{z_0}(q_0-1).
\]
For every nonzero $\det(H_I)$ the product over $j$ in \eqref{eq:sy-localizer} is consequently the same polynomial $\prod_j\det(T_j)^{Q_j-1}$. If $\det(H_{I_0})=0$, the initial factor $\Delta_0$ vanishes. Otherwise the right block of $H$ is nonzero, since the two planes differ. Cramer's rule gives a nonzero exchanged-column minor. The corresponding difference factor vanishes. The signs from ordering exchanged columns disappear because $Q_j-1$ is even. In the original chart each exchanged minor contains a normal variable, so only the $R_0$ term survives modulo $w$.
\end{proof}

We need slightly more of the last assertion. Expanding the product by nonempty subsets $S\subseteq\mathcal I$, each error term loses $|S|(Q_j-1)$ copies of $\Delta_j$. In the original chart, expanding its exchanged minors writes it as
\begin{equation}\label{eq:sy-local-error}
w^e X^{\nu}\prod_j(\Phi_g\Delta_j)^{\alpha_j-b_j}H,
\quad e\ge1,\quad |\nu|\le e,\quad 0\le b_j\le e,\quad H\in\mathcal T.
\end{equation}
Indeed, each exchanged determinant uses a normal column; after substitution a normal variable raised to $p^a$ contributes $w^{p^a}X^{p^a}$. Thus its normal degree is paid for by its $w$-power, and $e\ge\sum_j|S|(Q_j-1)\ge b_j$. This proves the source's more precise localizer expansion, not merely its reduction modulo $w$. \textsf{[SY2--SY3]}

\subsection{First-order keys and why top determinants occur}

For each position $s$, take the $r$ most recent even selected positions $a_l<s$, and let $T_s$ have entries $X_{a_l,b}^{p^{s-a_l}}$, $b<r$. Set $D_s=\det T_s$. A \emph{key} of order $d$ is an integral polynomial $K$ with a normal-homogeneous leading term $K_0$ of degree $d\le2Q$, such that
\[
K_0\in w^dO[X],\quad \Phi_gK-K_0\in w^{d+1}O[X],\quad
\Phi_{g'}K\in w^dO[X]\quad\hbox{for every chart }g'.
\]

For each normal variable $X_{s,b}$ there are order-one keys and tangential integral coefficients $H_l$ satisfying
\begin{equation}\label{eq:sy-key-identity}
wX_{s,b}D_s=\sum_lH_lK_{0,l}.
\end{equation}
Here is the construction. If $s$ is even, use the bordered $(r+1)$-minor with first row at $s$, the older $r$ rows raised to $p^{s-a_l}$, and columns $I_0\cup\{b\}$. Its leading term is $(-1)^rwX_{s,b}D_s$; all other terms contain $w^{p^{s-a_l}}\in w^2O$. If $s$ is odd, pair each older primal row, raised to that power, with the dual row at $s$. Their leading terms form $wT_s(X_{s,0},\ldots,X_{s,r-1})^{\mathsf t}$; multiplying by $\operatorname{adj}(T_s)$ proves \eqref{eq:sy-key-identity}. On any residue-plane chart the first keys vanish because $r+1$ vectors lie in an $r$-space, and the second because a vector and a covector annihilating that space pair to zero. This checks all key conditions. \textsf{[SY4]}

If a label monomial is $X^m$, the source's digit condition implies $m_{s,b}\le Q$. There are at most $kQ$ normal factors at each $s$. Multiplying \eqref{eq:sy-key-identity} for those factors therefore expresses
\[
w^{\deg_N m}X^m\prod_{s\in TT}D_s^{kQ}
\]
as a sum of products of key leading terms, of total key order $\deg_Nm$, with coefficients in $\mathcal T$. The unused determinant powers are nonnegative. This proves the residual-monomial assertion \texttt{RestMon}; no division by $D_s$ is used.

\subsection{Finite correction, with the exact determinant budget}

A term of level $\ell=i+\sum_a d_a$ has formal symbol and integral raw lift
\begin{align*}
S_{\mu,\lambda}&=w^i\prod_aK_{0,a}\,(\Phi_gR_0)^c
            \prod_j(\Phi_g\Delta_j)^\lambda\prod_sD_s^\mu H,\\
U_{\mu,\lambda}&=w^i\prod_aK_a\,R_{\rm loc}^c
            \prod_j\Delta_j^\lambda\prod_s\widehat D_s^\mu\widehat H.
\end{align*}
The hats denote integral polynomials before substitution, with $\Phi_g\widehat D_s=D_s$ and $\Phi_g\widehat H=H$. Thus $S_{\mu,\lambda}\in w^\ell O[X]$, while $\Phi_{g'}U_{\mu,\lambda}\in w^{\ell+c}O[X]$ on other planes.

Fix the desired precision $L=o+c$. Expanding $\Phi_gU-S$ and using \eqref{eq:sy-local-error} and \eqref{eq:sy-key-identity} gives, modulo $w^L$, finitely many new symbols at levels $\ell+E$, with
\begin{equation}\label{eq:sy-spares}
1\le E<L-\ell,\qquad
\mu'=\mu-(2Q+1)E,\qquad \lambda'=\lambda-E.
\end{equation}
For completeness, the coefficient estimate behind the budget is this: in a key error of order $d$, a monomial of normal degree $v$ has coefficient divisible by $w^x$, where $x=\max(d+1,v)$. Put $E=x-d$. Then $v\le(2Q+1)E$. Replacing its $v$ normal factors by \eqref{eq:sy-key-identity} uses at most that many copies of each $D_s$. A localizer error spends at most $E$ frame determinants. Products add both the gained levels and these costs; surplus determinants go into $H$. Terms already of level $L$ may be discarded. The coefficient bound uses integrality of $K$ before substitution, not just the displayed leading congruence. \textsf{[SY5]}

Starting at $\ell\ge o$ with $\mu=(2Q+3)c$, $\lambda=c$ leaves enough determinants at every step: the invariants are $\mu+(2Q+1)\ell$ and $\lambda+\ell$, and every correction has level at most $L-1$. Downward induction on $L-\ell$ now produces an integral $\widetilde U$ with
\[
\Phi_g\widetilde U\equiv S\pmod {w^{o+c}},\qquad
\Phi_{g'}\widetilde U\equiv0\pmod {w^{o+c}}\quad(\Pi(g')\ne\Pi(g)).
\]
Explicitly, replace each child symbol by its already corrected lift and subtract those lifts from $U$; the new error is in $w^L$. The process is finite, so it requires no completion or infinite limit. Projection to a fixed multidegree preserves all congruences because substitutions preserve positions.

\subsection{The actual symbol family and its remaining input}

The source's fixed multiplier is
\begin{equation}\label{eq:sy-sptrue}
\mathrm{SPtrue}=R_0^c\prod_j\Delta_j^c
                    \prod_{s\in TT}\widehat D_s^{(2Q+3)c+kQ}.
\end{equation}
Its substitution is integral and has normal degree zero. For an allowed path-generator polynomial $P$ of normal degree at most $T$, the two remaining generator assertions put its segment factors into the algebra generated by key symbols. The label factor is handled by SY4. Apply the finite correction to each normal-homogeneous component $\sigma_o$ of $P\Phi_g(\mathrm{SPtrue})$, $0\le o\le T$. It yields integral polynomials $F_o$ of the required multidegree satisfying
\begin{gather*}
\Phi_gF_o\equiv w^o\sigma_o\pmod {w^{o+c}},\qquad
\Phi_{g'}F_o\equiv0\pmod {w^{o+c}}\quad(\Pi(g')\ne\Pi(g)),\\
\sum_{o=0}^T\sigma_o=P\Phi_g(\mathrm{SPtrue}).
\end{gather*}
These are the actual separated realizations. The source's intermediate \texttt{Fs} are instead rational preimages of $\sigma_o$ under $\Phi_g$; distinguishing the two avoids an unwarranted assertion that those intermediate preimages are integral. \textsf{[SY6]}

Finally let $E'_t$ be the degree of \eqref{eq:sy-sptrue} at position $t$. Frame rows occur once per period, and at most $N_1+1$ top determinants involve a fixed position. Consequently
\begin{align*}
E'_t&\le\{c(1+rkQ)+c\}Q
 +(N_1+1)\{(2Q+3)c+kQ\}rQ\\
&\le15c(N_1+rk)Q^2=:E_0
\end{align*}
when $k\le c$, $N_1\le k$, $r\ge2$ and $Q\ge3$. Also $E'_t=0$ off $TT$. The source's size condition $a\ge Ck^3r^2NQ^5cY$, $C\ge30$, makes $E_0\le a$, so subtracting this degree budget from the prescribed degree is legitimate. This completes the realization and degree part of \texttt{SymbolRealisableBound}, conditional on the exact path-generator spanning and image assertions. The existence and dimension of the large selected symbol space, and the dual testing estimates, are separate obligations. \textsf{[SY7--SY8]}
